\documentclass[a4paper,10pt]{article}
\usepackage[T1]{fontenc}
\usepackage[utf8]{inputenc}
\usepackage{lmodern,textcomp}
\usepackage[margin=2.2cm]{geometry}
\usepackage{amsmath,amssymb,siunitx,booktabs,tabularx,array,enumitem,xcolor,fancyhdr,pifont}
\usepackage[colorlinks=true,linkcolor=blue!50!black,urlcolor=blue!50!black]{hyperref}
\setlength{\parindent}{0pt}
\setlength{\parskip}{4pt}
\setlength{\headheight}{14pt}
\renewcommand{\arraystretch}{1.2}
\renewcommand{\tabularxcolumn}[1]{>{\raggedright\arraybackslash}p{#1}}
\newcommand{\file}[1]{\nolinkurl{#1}}
\newcommand{\who}{\textcolor{gray}{[who: \rule{1.6cm}{0.4pt}]}}
\newcommand{\out}[1]{\par\hspace*{1.2em}\textit{Output:} #1}
\newcommand{\chk}[1]{\par\hspace*{1.2em}\textit{Check:} #1}
\newlist{todo}{itemize}{2}
\setlist[todo]{label=$\square$,leftmargin=1.6em,itemsep=5pt}
\sisetup{per-mode=symbol}
\pagestyle{fancy}\fancyhf{}
\fancyhead[L]{34870 Electroacoustics}
\fancyhead[R]{Loudspeaker project --- work plan}
\fancyfoot[L]{Group 10}
\fancyfoot[R]{\thepage}

\begin{document}
\begin{center}
{\Large\bfseries Loudspeaker project, System D: work plan}\\[5pt]
Group 10: Louis Andrianne, Sophie Kimura, Mads Rudolph\\[3pt]
Based on \file{34870_Project_Guide_2026.pdf}. Tick the boxes as we go, and put a name on each step.
\end{center}

\section*{How to use this}
The project follows the guide's five tasks. Each step below says what goes in (which file we
already have), what we do with it, what comes out, and how we check it before moving on. Do the
steps in order inside a task; Tasks~1 and~2 can run in parallel. Every result that ends up in
the quiz or the report gets a figure in \file{Project/figures/} and a line in
\file{Project/report/Project_report.tex}.

\begin{center}
\begin{tabularx}{\linewidth}{lX}
\toprule
Date & Milestone \\
\midrule
\textbf{Sun 8 Nov} & Project quiz part 1: LTspice models and the Lab D/E data in them (Tasks 1 and 2) \\
November & IEC listening room booked for all of November (Task 4) \\
\textbf{Mon 30 Nov} & Project quiz part 2: final results and discussion \\
\textbf{Thu 3 Dec} & Demonstration in room 019: bring the passive board and the DSP coefficients \\
\bottomrule
\end{tabularx}
\end{center}

\section*{What we already have}
\begin{tabularx}{\linewidth}{lX}
\toprule
Data & File (labs repo) \\
\midrule
Woofer impedance, free air and closed box & \file{Lab D/data/labD_woofer_free.mat}, \file{labD_woofer_closed.mat} \\
Woofer impedance, one vent at 160/200/240 mm & \file{Lab D/data/labD_vent_L160/L200/L240.mat} \\
Near field of cone and vent, each length & \file{Lab D/data/labD_nf_cone_L*.mat}, \file{labD_nf_vent_L*.mat} \\
Midrange and tweeter impedance $Z_i$ & \file{Lab D/data/labD_midrange.mat}, \file{labD_tweeter.mat} \\
Anechoic $p/V$ of all three units, $0^\circ$ and $30^\circ$ & \file{Lab E/data/labE_system_0.mat}, \file{labE_system_30.mat} \\
Woofer data sheet & \file{Lab D/datasheets/} \\
Course tools & \file{Project/matlab/course/} (\file{Matlab2LTspice.m}, \file{DigitalCrossover.m}, \file{ventbox}) \\
Course LTspice examples & \file{Project/ltspice/course/passive/}, \file{digital/} \\
One loader for all of it & \file{Project/matlab/load_systemD.m} \\
\midrule
\multicolumn{2}{l}{\textbf{Ready to use (made from the data above)}} \\
LTspice data files, names already unique & \file{Project/data/Zi_mid.txt}, \file{Zi_tw.txt}, \file{TF_\{woof,mid,tw\}.txt},
  \file{TF30_*.txt} (made by \file{matlab/export_systemD.m}, $d = \SI{-1.15}{m}$) \\
The same data as CSV (Python, Excel) & \file{Project/data/csv/} \\
One test schematic per unit & \file{Project/ltspice/units/Unit_woof.asc}, \file{Unit_mid.asc}, \file{Unit_tw.asc} \\
Empty system skeleton (Figure~1) & \file{Project/ltspice/System_template.asc} \\
\bottomrule
\end{tabularx}

The unit and system files are written by \file{Project/ltspice/gen_templates.py};
\file{python3 gen_templates.py --verify} runs LTspice and confirms each unit reproduces the
measured $Z_i$ and $p/V$ (to \SI{0.04}{dB}). Open them in LTspice straight from the repo:
the \file{.inc} paths are relative.

Known problems with the data (see \file{Project/README.md}): $Z$ already uses $R = \SI{32.9}{\ohm}$
and $V_\text{amp} = 2\cdot\text{AI0}$; temperature was not recorded (assume \SI{20}{\celsius});
the UMIK channel is unusable above \SI{18.4}{kHz}; \file{labE_system_30_bad.mat} is rejected;
the woofer TF in Lab E was taken with the vent plugged.

\section*{Task 0: Before we start}
\begin{todo}
\item Agree who owns which task, and book lab time with Group 4 (we share System D). \who
\item Measure what we never wrote down, next time we are at the box: inner diameter of
  \emph{both} vents, how the telescopic length is defined (and its min/max), whether there was
  filling in the box during the Lab D runs, and the room temperature. \who
  \chk{Task 1 needs the second vent's diameter to model ``two vents''. Without it we can only
  model the one we measured.}
\item Everyone runs \file{S = load_systemD} in MATLAB once, so we know the data loads on every PC. \who
\end{todo}

\section*{Task 1: Low-frequency tuning of the woofer box}
Lectures 7 and 8, Lab D. Goal: choose closed or vented, the vent mass (none, one, two vents) and
filling, from an LTspice analogy that has been validated against our near-field measurements.

\subsection*{1.1 Thiele--Small parameters}
\begin{todo}
\item Run \file{Lab D/matlab/analyse_labD.m} and copy $f_S$, $Q_{MS}$, $Q_{ES}$, $Q_{TS}$,
  $V_{AS}$, $R_E$ into a table next to the data-sheet values. \who
  \out{T-S table for the report (section 2.1).}
  \chk{$f_c/f_S$ should equal $Q_{TC}/Q_{TS}$ (the script prints both). Explain why our
  $V_{AS}$ differs from the data sheet.}
\item Derive the circuit elements $M_{MS}$, $C_{MS}$, $R_{MS}$, $Bl$ from them ($S_D$ from the
  data sheet), and decide how much of $M_{MS}$ is air load (Lecture 7: $M_{MD} = M_{MS} - $ air
  mass). \who
\end{todo}

\subsection*{1.2 LTspice analogy circuit}
\begin{todo}
\item Draw the woofer-in-box circuit in LTspice (Lecture 7 slide 27 / Lecture 8): electrical,
  mechanical and acoustical domain, front radiation impedance, box compliance $C_{AB}$. Start
  from the closed box. \who
  \chk{Simulated $|Z_E|$ against \file{labD_woofer_closed.mat}: same peak frequency and height.}
\item Add the vent: acoustic mass $M_{AP} = \rho(L + \Delta L)/S_P$ in parallel with $C_{AB}$,
  plus a leak/loss resistance. Write down which end correction we use and why. \who
  \chk{Simulated $|Z_E|$ against \file{labD_vent_L160/200/240.mat}: both peaks and the minimum
  between them ($\approx f_B$). If the formula does not move $f_B$ the way the measurement does,
  write that down; it tells us something about the vent geometry.}
\item Add outputs: far-field pressure $\propto j\omega(U_D - U_P)$ and the near-field pressures
  of cone and vent ($\propto j\omega U/a$). \who
\end{todo}

\subsection*{1.3 Validation against the near field}
\begin{todo}
\item Plot simulated against measured near field (cone and vent) for each vent length, using
  \file{labD_nf_*}. \who
  \out{One figure per length, or one figure with three panels.}
  \chk{Cone notch at $f_B$, vent peak around $f_B$, crossing point of the two curves. Explain any
  offset (mic position, room modes below \SI{120}{Hz}, the vent mic picking up the cone above
  $\approx$\SI{60}{Hz}).}
\end{todo}

\subsection*{1.4 Choose the tuning}
\begin{todo}
\item With the validated circuit, step through the options the guide allows: closed; one vent at
  different lengths; two vents in parallel ($M_{AP,1} \parallel M_{AP,2}$); filling (larger
  effective volume, more damping). \who
\item For each one, read off $f_3$ and the ripple of the low-frequency response. \who
  \out{A table option $\to$ $f_B$, $f_3$, ripple, and one overlay plot.}
\item Pick one and argue for it (bandwidth against ripple, cone excursion near $f_B$, what the
  standard alignments in Lecture 8 would ask for, and why we cannot reach them). \who
\end{todo}

\section*{Task 2: Modelling the crossover filters}
Lecture 9, Lab D and Lab E.

\subsection*{2.1 Measured data into LTspice}
\begin{todo}
\item Read \file{matlab/export_systemD.m} and understand what it does: it calls
  \file{Matlab2LTspice.m} for $Z_i$ of midrange and tweeter (no distance) and the TF of all three
  units (same $d$ for all), then renames the first line of each file (\file{Vspkr_mid},
  \file{Pout_mid}, \dots) because duplicate names break the simulation. Try another $d$ if the
  phase is hard to read, and re-run it. \who
\item Open \file{ltspice/units/Unit_mid.asc}, \file{Unit_tw.asc}, \file{Unit_woof.asc} and run
  them. \who
  \chk{LTspice $|Z|$ and phase equal Lab D; \file{V(Pout_*)} equals Lab E's $p/V$ (the script
  checks this, but look at the plots yourselves and put one in the report).}
\item Note: the course's \file{...and example filter.asc} is drawn as a \textbf{high-pass}
  (C in series, L to ground) although the guide's Figure~2 calls it a low-pass. \who
\end{todo}

\subsection*{2.2 Choose the two crossover frequencies}
\begin{todo}
\item Overlay the three measured responses (Lab E, $0^\circ$ and $30^\circ$) and the impedances. \who
\item Woofer--midrange: where both are still smooth, and where the woofer has not started to
  beam (compare $0^\circ$ and $30^\circ$). \who
\item Midrange--tweeter: well above the tweeter resonance (Lab D), and where the midrange still
  behaves off axis. Note the level difference between the two units. \who
  \out{Two frequencies with one line of justification each.}
\end{todo}

\subsection*{2.3 Digital crossover (woofer / mid-high)}
\begin{todo}
\item Run \file{DigitalCrossover.m} for the LPF and the HPF at our frequency; decide the order. \who
  \chk{The function makes Butterworth sections ($n=4$ gives $Q = 0.541$ and $1.307$). If we
  want Linkwitz--Riley, cascade two $n=2$ sections at the same $f_0$, and say so in the report.}
\item Paste the Laplace strings into the E sources of \file{ltspice/course/digital/}, woofer branch
  first, and check the filter on its own. \who
\end{todo}

\subsection*{2.4 Passive crossover (midrange / tweeter)}
\begin{todo}
\item Choose topology and order; compute the values for the nominal impedance. \who
\item Round to the standard values in Appendix 1 (no combining components). \who
\item Simulate each filter loaded with a resistor, then with the real $Z_i$, and compare. \who
  \chk{Does the driver impedance shift the corner? Add a Zobel or a series resistor if it does.}
\item Match the tweeter level to the midrange (L-pad or series resistor). \who
\item Tolerance check: step the component values by their tolerance and look at the sum. \who
\end{todo}

\subsection*{2.5 Complete system}
\begin{todo}
\item Start from \file{ltspice/System_template.asc}: it is Figure~1 already wired, with gain-1
  biquad slots and plain wires where the passive filters go. Then wire all three units as in the guide's Figure 1 (digital LPF $\to$ woofer; digital HPF
  $\to$ passive LPF/HPF $\to$ mid/tweeter), and sum the three outputs with series controlled
  sources. Try polarity $\pm1$ per unit. \who
\item Add the woofer/mid level difference as a gain (this becomes the DSP output gain). \who
\item Iterate (frequencies, orders, polarity, levels) until the on-axis sum is flat and the
  $30^\circ$ sum is reasonable. Keep a log of each iteration and why it changed. \who
  \out{Final LTspice schematic, on-axis and $30^\circ$ sum, iteration log.}
\end{todo}

\subsection*{Quiz part 1 (Sun 8 Nov)}
\begin{todo}
\item Everyone answers individually, but compare figures beforehand: T-S table, validated box
  circuit, near-field comparison, LTspice $Z_i$/TF checks, first full-system sum. \who
\end{todo}

\section*{Task 3: Implementing the crossovers}
\begin{todo}
\item MiniDSP: paste the \file{biquad_DSP} strings into Xover of outputs 1 and 2 (everything
  else bypassed); set the output gains from 2.5. \who
\item Passive board: mount our standard values; photograph it; note which parts we used. \who
\item Anechoic measurement of the complete system (vents may be covered), on and off axis. \who
  \chk{Measured against simulated. Explain the differences and decide if another iteration of
  Task 2 is needed.}
\end{todo}

\section*{Task 4: Listening test (IEC room, November)}
\begin{todo}
\item Read Lecture 10 (and Bech \& Zacharov if needed) and write a short procedure: what we
  compare, in what order, how long, listening position, level. \who
\item Choose the programme material (a few tracks with clear bass, voice, cymbals). \who
\item Write the questionnaire (fixed attributes and scale), and find subjects. \who
\item Run the test, collect the answers, summarise them, and decide what to change. \who
\end{todo}

\section*{Task 5: Quiz part 2 and demonstration}
\begin{todo}
\item Final results and discussion in quiz part 2 (Mon 30 Nov), individually. \who
\item Demonstration Thu 3 Dec: passive board and DSP coefficients ready, plus the plots we
  want to show. \who
\item Fill in \file{Project_report.tex} section by section as each task finishes. \who
\end{todo}

\end{document}
